Put
\[
 L:=\log(en),\qquad b:=\frac1n+\frac d{n^2},\qquad
 u:=\frac{d^2}{n^2L^2},\qquad m:=\min\{u,\sigma^2\}.
\]
In the residual wedge, \(d/n^2<u<1\).  The elementary minimum comparison in Theorem~\(\mathrm{thm:fixed\text{-}interior\text{-}tightness\text{-}and\text{-}shrinking\text{-}radius\text{-}gap}\) gives
\[
 \frac12\{b+m\}
 \le \frac1n+\min\left\{u,\sigma^2+\frac d{n^2}\right\}
 \le b+m,
\]
whereas Theorem~\(\mathrm{thm:radius\text{-}channel\text{-}converse}\) gives the normalized lower benchmark \(b+\sigma^2u\).  Thus only the stated joint wedge remains.

There is an exact lower-side radius parametrization.  For a pair having total mass \(\omega>0\) and share \(r\in[0,1]\), set
\[
 p_+:=\omega r,\qquad p_-:=\omega(1-r),
 \qquad
 \vartheta_+:=\theta+\sigma(1-r),\qquad
 \vartheta_-:=\theta-\sigma r.
\]
Then
\[
 \frac{p_+\vartheta_++p_-\vartheta_-}{p_++p_-}=\theta,
 \qquad
 |\vartheta_+-\theta|\le\sigma,\qquad
 |\vartheta_--\theta|\le\sigma.
\]
Moreover, when \(\sigma>0\), the translation
\[
 \theta'=\theta+\Delta,\qquad r'=r+\Delta/\sigma
\]
leaves both cell effects unchanged whenever \(r'\in[0,1]\), while changing the pair-average effect by \(\Delta\).  Hence a separation of order \(\min\{\sigma,d/(nL)\}\) is algebraically compatible with realization-wise radius control.  If \(\omega\) is deterministic, however, equality of the first observed intensity moment for the plus cell forces
\[
 \omega E_0[r]=E_0[p_+]=E_1[p_+]=\omega E_1[r],
\]
which rules out a nonzero mean translation of this form.  The remaining lower-side task is therefore a correlated moment problem for \((\omega,r)\), together with normalization concentration and fixed-sample de-Poissonization; it is not supplied by either frozen binary lower-bound lemma.

The natural equal-mass reflected pairing has an additional obstruction already at the second observable moment.  Give the two cells common raw mass \(q\), normalized effects \(t_h+e\) and \(t_h-e\), and propensities \(\pi_+\) and \(\pi_-\), where \(|e|\le\sigma\).  Write \(v_\pm:=\pi_\pm(1-\pi_\pm)\).  By the definition of the collision weights, the pair's expected collision numerator, apart from the positive factor \(2q^2M\), is
\[
 (v_++v_-)t_h+(v_+-v_-)e.
\tag{1}
\]
For the usual reflected choice \(\pi_-=1-\pi_+\), one has \(v_-=v_+\), so \((1)\) directly reveals every nonzero shift of \(t_h\).  Symmetric reflection therefore cannot satisfy even the required degree-two moment matching.  An asymmetric pair can cancel this one moment only by coupling its deviation and propensities; all higher mixed moments remain to be matched jointly.

For the upper-side reduction, normalize
\[
 V:=Y/M,\qquad \vartheta:=\tau(P)/M,\qquad
 \vartheta_k:=\tau_k/M,\qquad
 s_{ak}:=P(X=k,A=a),\qquad
 z_{ak}:=E_P[V\mathbf 1\{X=k,A=a\}].
\]
For an independent pilot value \(t\in[-1,1]\), define
\[
 H_k(t):=s_{0k}z_{1k}-s_{1k}z_{0k}-t s_{0k}s_{1k}
 =s_{0k}s_{1k}(\vartheta_k-t).
\]
Let \(r_K\) be a reciprocal polynomial on the light-cell range, extend it by \(r_K(x):=x^{-1}\) on the population heavy-cell range, and put
\[
 e_{ak}:=1-s_{ak}r_K(s_{ak}),\qquad
 c_k:=s_{0k}s_{1k}\{r_K(s_{0k})+r_K(s_{1k})\},
 \qquad
 S_K(t):=\sum_kH_k(t)\{r_K(s_{0k})+r_K(s_{1k})\}.
\]
Since \(p_k=s_{0k}+s_{1k}\), direct algebra gives
\[
 p_k-c_k=s_{1k}e_{0k}+s_{0k}e_{1k}.
\tag{2}
\]
Using \(\sum_kp_k(\vartheta_k-\vartheta)=0\), another direct expansion yields the exact identity
\[
 t+S_K(t)-\vartheta
 =\sum_k(\vartheta_k-\vartheta)(c_k-p_k)
 +(t-\vartheta)\left(1-\sum_kc_k\right).
\tag{3}
\]
Suppose the polynomial satisfies
\[
 |e_{ak}|\le \frac{B}{K^2s_{ak}}
\tag{4}
\]
on the light range.  Assumption~\(\mathrm{ass:overlap}\) implies that \(s_{1k}/s_{0k}\) and \(s_{0k}/s_{1k}\) are bounded by constants depending only on \(\epsilon\).  Equations~\((2)\)--\((4)\), together with Assumption~\(\mathrm{ass:approximate\text{-}homogeneity}\), therefore give
\[
 |t+S_K(t)-\vartheta|
 \le C_{\epsilon}\{\sigma+|t-\vartheta|\}\frac{dB}{K^2}.
\tag{5}
\]
Taking \(B\asymp L/n\) and \(K\asymp L\) makes the square of the right side in \((5)\) at most
\[
 C_{\epsilon}u\{\sigma^2+(t-\vartheta)^2\}.
\tag{6}
\]

This proves a conditional reduction of the upper route.  If a total clipped heavy--light factorial lift \(\widehat S_K(t)\) of \(S_K(t)\) were shown to satisfy
\[
 E[(\widehat S_K(t)-S_K(t))^2\mid t]
 \le C_{\epsilon}\left\{b+u[\sigma^2+(t-\vartheta)^2]\right\},
\tag{7}
\]
then an independent collision pilot from Lemma~\(\mathrm{lem:continuous\text{-}occupancy\text{-}collision\text{-}upper}\) would obey
\[
 E[(t-\vartheta)^2]\le C_{\epsilon}(b+\sigma^2).
\tag{8}
\]
Clipping \(t+\widehat S_K(t)\) to \([-1,1]\) cannot increase squared error because \(\vartheta\in[-1,1]\).  Combining \((5)\)--\((8)\) and using \(u<1\) would give normalized risk \(C_{\epsilon}(b+\sigma^2u)\), exactly the proved product converse benchmark.  The unresolved statement is \((7)\): the existing linear-mark covariance lemma neither contains its centered mixed-arm kernel nor proves its exact-homogeneity degeneracy.  No inference claim and no adaptation to unknown \(M\) or \(\sigma\) is used.